\section{Conclusion}

This project demonstrated the potential of combining audio and lyrical information for music understanding and recommendation, while also highlighting several practical and conceptual limitations. In particular, inconsistent alignment between lyrical meaning and musical character, the manual construction of user profiles, infrastructure-related obstacles, and incomplete lyric coverage all constrained the scope of the evaluation.

Given more time, a useful extension would be to automate preprocessing by writing a script that segments longer tracks into 30-second clips and averages the resulting predictions. Another promising direction would be to develop a representation that captures the semantic content conveyed by musical notes themselves. For example, one could draw on ideas from music theory to group chords, motifs, or harmonic progressions and model them as a “corpus” analogous to natural language, thereby enabling analysis even for tracks without lyrics.